% Chapter 20, typeset from lectures/L20/README.md and info/examination.md.

\chapter{Practical exam 2: CAN modules}
\label{c20:ch}

{\large\itshape Three hours, individual, at the computer. Covers the project's modules: bit timing,
CRC, bit stuffing, frame sequencing, register semantics and the SPI transaction.\par}

\begin{chapterbox}{In this chapter}
\begin{itemize}
  \item Why an individual exam follows a group project.
  \item The exam's format: writing tasks with a testbench handed out, and reading tasks on paper.
  \item The four kinds of task: variants, repairs, a single register, reading.
  \item What you are allowed to use \textemdash{} including the group's own project code.
  \item How to prepare, what the exam measures, and where the design goes afterwards.
\end{itemize}
\end{chapterbox}

This is the course's second and last examination. Like \chapref{c9:ch} it contains no new theory; the
chapter says what the exam measures, what shape the tasks take and how to prepare. Marks and grade
boundaries are in appendix~\ref{app:project}.

The exam covers \textbf{the project's modules}, that is, \chapref{c10:ch} to \ref{c19:ch}: the bit
timing, the CRC, the bit stuffing in both directions, the frame sequencing, the register semantics
and the SPI transaction. Part~\ref{part:vhdl}'s material has not gone away \textemdash{} every module
in the project is built from the synchronous process template, a counter, a shift register and a
Moore machine \textemdash{} but the emphasis is on part~\ref{part:can}.

% ------------------------------------------------------------------------------------------
\section{Why an individual exam on a group project}
\label{c20:sec:why}

The project is assessed as a group and accounts for half the course. This exam measures something
else: that every individual understands what the group built, not only the part they wrote
themselves.

That is not a distrust of the way of working but a consequence of it. A group of four or five splits
eleven modules between them, and that split is exactly what makes the project feasible in ten
sessions. The downside is equally predictable: whoever wrote \code{crc15} can get away without ever
having needed to understand why \code{bit_timer} has two pulses, and whoever wrote the register bank
without having thought about what the stuffing does to the CRC. The tasks are therefore drawn from
\emph{all} parts of the project, regardless of who wrote which module.

\begin{aside}
There is a practical consequence of this worth taking seriously long before the exam: \textbf{the code
review during the project is the preparation}. A group that genuinely reviews each other's PRs,
instead of approving them unread, has already read everything the exam asks about.
\Chapref{c15:ch}'s code review seminar is the occasion to practise that; it is also the occasion that
pays off here.
\end{aside}

% ------------------------------------------------------------------------------------------
\section{Format}
\label{c20:sec:format}

The same shape as \chapref{c9:ch}: independent tasks, unrelated to each other and ordered roughly by
increasing difficulty.

Most are \textbf{writing tasks} \textemdash{} a self-checking testbench handed out and a description
of the module it drives, with port order, types and behaviour. The task is to write the module so that
the testbench passes.

Besides those there are \textbf{reading tasks}: a timing diagram or an SPI transaction to interpret on
paper. They have no testbench, for exactly the reason that makes them worth asking \textemdash{}
reading a waveform is a skill in itself, and it is not examined by writing VHDL. They are marked
against an answer key, and partly correct answers earn partial marks the same way partly correct logic
does.

\subsection*{The four kinds of task}
\begin{description}[style=unboxed,leftmargin=0pt,itemsep=0.8ex]
  \item[Variants.] A \code{bit_timer} with a different sample point or a different bit rate. A
    stuffing detector counting a different number. A CRC with a different polynomial. The module is
    familiar; the one point where it differs from the project's is the point that requires you to
    have understood it.
  \item[Repairs.] A given state machine or a given shift register that is \emph{almost} right. Find
    the fault using the testbench's output and fix it. It is \chapref{c14:ch}'s and \ref{c15:ch}'s
    working method put to the test: read what the testbench reports, and work backwards from there to
    the line that caused it.
  \item[A single register.] One register from the register bank, with its set and clear rule, isolated
    from the rest. \Chapref{c18:ch}'s difference between a pulse and a sticky level is the whole
    task.
  \item[Reading.] A timing diagram or an SPI transaction to interpret: what is on the bus, and what
    should have been there?
\end{description}

Partly correct logic earns partial marks, so never hand in an empty file because the module did not
get finished; the same goes for a half-interpreted timing diagram. The marking reads the code as well,
not only the testbench's outcome: a module that passes by doing something other than what was asked
does not earn full marks.

% ------------------------------------------------------------------------------------------
\section{Permitted aids}
\label{c20:sec:aids}

\textbf{Allowed:}
\begin{itemize}
  \item All the course material: this book, the lectures' READMEs and every appendix. The register map
    and the protocol specification (appendix~\ref{app:protocol}) belong there of course; several tasks
    are written against them.
  \item \textbf{The group's own project code.} You built it; it would be strange to pretend
    otherwise. The tasks are written so that copying does not help \textemdash{} the variants differ
    at precisely the point that requires you to have understood the module, and the repairs have no
    counterpart in the repository at all.
  \item Your own computer with GHDL, and CircuitVerse if you want to draw a circuit or a waveform
    before you write it. That is often faster, not slower.
\end{itemize}

\textbf{Not allowed:} communication with others in any form; language models and other tools that
generate or explain VHDL for you; and searching the internet for ready-made solutions. The line is not
drawn at who presses the keys but at whose understanding is being tested.

The full and authoritative list of permitted aids is in appendix~\ref{app:project}; if anything here
says otherwise, that one applies. If you are unsure whether something is allowed: ask before the exam,
not during it.

% ------------------------------------------------------------------------------------------
\section{Before the exam}
\label{c20:sec:prep}

\begin{itemize}
  \item Go through the modules in the project that \emph{somebody else} in the group wrote. That is
    where the gaps are, and it is the one piece of advice in this list that is hard to follow the day
    before.
  \item The self-check questions at the end of \chapref{c12:ch} to \ref{c19:ch} are deliberately
    written as revision material. If you can answer them all without looking anything up, you are
    ready.
  \item Read the testbenches handed out through once more. They are the clearest thing there is about
    exactly what every module promises, and a writing task is always easier for somebody who knows how
    a testbench states its requirements.
  \item Redraw \code{bit_timer}'s two pulses and \code{can_controller}'s state diagram on a sheet of
    paper, from memory. Those two pictures carry more of the project than any single file does.
  \item Make sure GHDL works on your machine \emph{before} you arrive. A broken toolchain is no reason
    for extra time.
\end{itemize}

% ------------------------------------------------------------------------------------------
\section{What the exam measures}
\label{c20:sec:measures}

\begin{itemize}
  \item That you can turn a specification in words into a correct module, with the right port order.
  \item That you understand the bit timing: why the sample point sits where it does, and which pulse
    does what.
  \item That you can reason about bit stuffing in both directions \textemdash{} inserting and removing
    are the same rule seen from two sides, and both sides have to keep the CRC out of it.
  \item That you can tell a pulse from a sticky level, and know which layer is responsible for which.
  \item That you can read a testbench's output and find your way back to the cause.
\end{itemize}

% ------------------------------------------------------------------------------------------
\section{After the course}
\label{c20:sec:after}

The controller goes on: the parallel class writes the C++ driver against the same registers over the
same SPI transport, and the node travels on to the CAN labs in the next embedded systems course. What
holds the two halves together is the register map and the protocol specification in
appendix~\ref{app:protocol} \textemdash{} not the code on either side. That is the whole point of an
interface: two groups that have never read each other's source can still meet on a bus, because the
specification between them was written before either of them started.

Over twenty chapters the book has gone from a gate on paper to a CAN node a microcontroller can talk
to. Every pattern along the way is common property in digital design and carries over unchanged into
whatever you build next:

\begin{itemize}
  \item the two-flip-flop synchroniser (\chapref{c4:ch}), which every asynchronous input passes
    through,
  \item the clocked process template (\chapref{c3:ch}), which every sequential module in the book is
    written with,
  \item the Moore machine (\chapref{c8:ch}), as an enumerated type and a \code{case},
  \item the shift register and the counter (\chapref{c6:ch}) as building blocks rather than as special
    cases, and
  \item finally the split into blocks with one clear responsibility each (\chapref{c11:ch}), which is
    the only reason eleven modules can be written by five people in parallel and still fit together.
\end{itemize}

If you want to go further than the design in this book \textemdash{} to a controller with error
counters and ACK checking, to verification beyond a hand-written testbench, or to the next protocol
built with the same method \textemdash{} appendix~\ref{app:reading} points out what there is to read
next.
